\documentclass[11pt,a4paper]{article}
\usepackage[utf8]{inputenc}
\usepackage{amsmath,amssymb,amsthm}
\usepackage{booktabs}
\usepackage{hyperref}
\usepackage{geometry}
\usepackage{pgfplots}
\usepackage{orcidlink}
\pgfplotsset{compat=1.18}
\geometry{margin=1in}

\newtheorem{theorem}{Theorem}
\newtheorem{conjecture}{Conjecture}
\newtheorem{remark}{Remark}

\title{Empirical Falsification of Convergence Rate Predictions\\
for the Connes--Consani--Moscovici Zeta Spectral Triple}

\author{Ronnie Andrews, Jr.\,\orcidlink{0009-0003-9724-3104}\\
\textit{Team Xcelerator Inc.\textsuperscript{\textregistered}}\\
\texttt{randrewsmath@gmail.com}}

\date{June 2026}

\begin{document}
\maketitle

\begin{abstract}
The Connes--Consani--Moscovici (CCM) operator construction
(arXiv:2511.22755) produces eigenvalues that match imaginary parts
of nontrivial Riemann zeta zeros to remarkable precision (55~digits
at $(\lambda^2{=}13, N{=}120, 200\text{-digit})$; up to 999~digits
at $\lambda^2{=}1000$ in our independent extension). Two papers in the surrounding
literature make quantitative predictions about the convergence
rate of this construction: CCM Lemma~7.2 conjectures that the
smallest Weil-form eigenvector $\xi_\lambda$ is approximated by a
prolate-wave educated guess $k_\lambda$ at relative rate
$O(\lambda^{-2})$, and \'Sliwi\'nski's Conjecture~4.1 (arXiv:2601.12133)
predicts $\varepsilon(\kappa)\cdot\ln(\kappa)$ approaches a nonzero
constant in the regime $\kappa = N = \lambda$. We test both
predictions at high precision (primarily $1000$ digits via
MPFR/GMP).
Conjecture~7.2 is empirically falsified: $\text{rel}\times\lambda^2$
grows by approximately $1.116\times 10^6$ across $\lambda^2 \in
[13, 1000]$ rather than remaining bounded by a constant.
Conjecture~4.1 is empirically unsupported at HP-1000: across
$\kappa = N = \lambda \in [50, 500]$,
$\varepsilon\cdot\ln\lambda$ ranges over $[7.792, 32.09]$
(full-spectrum, $4.118\times$ spread) and $[1.250, 10.56]$
(trimmed-interior, $8.445\times$ first/last), without
approaching a constant, with a sharp upturn at $\kappa = 500$.
Theorem~3.1 (\'Sliwi\'nski's lower bound) is rigorously satisfied
at every configuration with substantial margin
($5.001\times$--$128.4\times$ ratios). Per-eigenvalue
analysis reveals that the construction converges
\emph{super-exponentially} on the first eigenvalues --- the first
Riemann zero is matched to $160.8$~digits at $\kappa = 50$ and
to $608.4$~digits at $\kappa = 200$, with the matching-digit
count growing roughly linearly in $\kappa$ before saturating
against the $1000$-digit working-precision ceiling at $\kappa
= 400$ (accuracy is controlled jointly by $\lambda$, $N$, and
working precision; the regime $\kappa = N = \lambda$ scales all
three together). The ``slow convergence''
measured by aggregate error metrics is dominated by high-index
eigenvalues that lie at the matrix's truncation boundary. We
further show that the CCM
construction is $10^{55}$ more accurate than naive Mellin
truncation (verified at HP-1000), confirming its power is
algebraic rather than analytic. The correct convergence question
is not ``what is the rate?'' but ``how many eigenvalues converge
to full precision as $\kappa$ grows?''
\end{abstract}

\section{Introduction}

The construction of Connes, Consani, and Moscovici~\cite{CCM2025}
proposes a finite-dimensional operator family
$D_{\log}^{(\lambda, N)}$ whose spectrum approximates the
imaginary parts of the nontrivial Riemann zeros with extraordinary
numerical accuracy. The construction has been independently
reproduced and extended to 999~matching digits at
$(\lambda^2{=}1000, N{=}800, \text{HP-1000})$ in our companion
paper~\cite{Andrews2026A}.

A natural question is the rate at which this approximation
improves with $\lambda$ and $N$. Two recent works in the
literature offer specific quantitative predictions:

\begin{conjecture}[CCM Lemma~7.2~\cite{CCM2025}]
\label{conj:lemma72}
There exists a constant $C > 0$ such that the smallest Weil-form
eigenvector $\xi_\lambda$ is approximated by the prolate-wave
educated guess $k_\lambda$ in relative $L^\infty$ error at rate
$\|\xi_\lambda - c\,k_\lambda\|_\infty / \|\xi_\lambda\|_\infty
\leq C\,\lambda^{-2}$ for some optimal scalar~$c$.
\end{conjecture}

\begin{conjecture}[\'Sliwi\'nski Conjecture~4.1~\cite{Sliwinski2026}]
\label{conj:sliwinski}
In the regime $\kappa = N = \lambda$, the limit
$\lim_{\kappa\to\infty} \varepsilon(\kappa)\cdot\ln(\kappa)$
exists and equals a nonzero constant, where
$\varepsilon(\lambda, N) = (1/N)\sum_{k=1}^{N}|\nu_k - \zeta_k|$
is the mean absolute spectral error.
\end{conjecture}

Both conjectures, if true, would provide a route to proving the
Riemann Hypothesis: Conjecture~\ref{conj:lemma72} is Step~2 of
the proof strategy in~\cite{CCM2025}, and
Conjecture~\ref{conj:sliwinski} would imply $\varepsilon \to 0$
hence eigenvalue convergence to the Riemann zeros.

In this paper we:
\begin{enumerate}
\item Test Conjecture~\ref{conj:lemma72} at high precision
(primarily $1000$ digits) across $\lambda^2 \in \{13, 100, 1000\}$ and
find it empirically falsified: the quantity
$\text{rel}\times\lambda^2$ grows by a factor of approximately
$1.116\times 10^6$ across the tested range;
\item Test Conjecture~\ref{conj:sliwinski} in its own regime
($\kappa = N = \lambda$) at HP-1000 ($3338$-bit precision) across
$\kappa \in \{50, 100, 150, 200, 300, 400, 500\}$.
$\varepsilon(\kappa)\cdot\ln\lambda$ does not approach a constant:
full-spectrum values range over $[7.792, 32.09]$ ($4.118\times$
spread), trimmed-interior values over $[1.250, 10.56]$
($8.445\times$ first/last), with a sharp upturn at $\kappa = 500$.
The conjectural limit form is empirically unsupported by this
data. Theorem~3.1 (\'Sliwi\'nski's lower bound) is satisfied at
every configuration with substantial margin
($5.001\times$--$128.4\times$ ratios);
\item Quantify how much of CCM's accuracy is captured by integral
transforms of $\xi_\lambda$, at HP-1000 (3338-bit) precision: at
$\lambda^2 = 13$, naive Mellin truncation gives $O(1)$ error
against Riemann zeros while CCM gives $55$~matching digits, an
improvement factor of $10^{55}$;
\item Argue from this evidence that integral-transform bridges
from $\xi_\lambda$ to the Riemann zeros (Mellin truncation and
$\xi$-weighted variants tested here) fall short of CCM by orders
of magnitude, and that the construction's accuracy comes from
\emph{algebraic} structure
(eigenvalue problem followed by rational-function zeros).
\end{enumerate}

The CCM construction itself is not in question here. Our
companion paper~\cite{Andrews2026A} verifies its eigenvalues
against $2000$-digit reference zeros (Step~1 of the CCM proof
strategy holds to $999$~matching digits at $\lambda^2{=}1000$,
$N{=}800$, HP-1000). The findings of the
present paper concern only the proposed quantitative bridges
\emph{from} $\xi_\lambda$ \emph{to} the Riemann Hypothesis, and
the accuracy comparison against alternative integral
formulations. All computations are reproducible from the
source-available implementation at\\
\url{https://github.com/TeamXcelerator/ccm-convergence-rate-falsifications}.

\section{Mathematical Setup}

We briefly recall the relevant constructions, following the
notation of~\cite{CCM2025} and~\cite{Sliwinski2026}.

\subsection{The Weil Quadratic Form and the Eigenvector $\xi_\lambda$}

Let $\lambda > 1$ and $L = 2\ln\lambda$. The Weil quadratic form
$QW_\lambda$ is a real symmetric $(2N{+}1) \times (2N{+}1)$ matrix
with entries
\[
\tau_{n,m} = W_{0,2}(V_n, V_m) - W_R(V_n, V_m)
- \sum_{\substack{k = p^j \\ k \leq \lambda^2}}
\frac{\log p}{\sqrt{k}}\, q(V_n, V_m)(\log k),
\]
where $V_n(u) = (1/\sqrt{L})\,e^{2\pi i n \ln(\lambda u)/L}$
forms an orthonormal basis on
$\mathcal{H}_\lambda = L^2([\lambda^{-1}, \lambda], du/u)$.
The smallest eigenvalue $\varepsilon_N$ of $\tau$ has eigenvector
$\xi = (\xi_{-N}, \ldots, \xi_N)$, symmetrized to be even
($\xi_{-j} = \xi_j$) and normalized so $\sum_j \xi_j = \sqrt{L}$.
(The forced-even symmetrization is a theoretical guarantee; empirically
the natural smallest eigenvector is even without it at every
configuration tested above the precision floor --- see~\cite{Andrews2026A}.)

\subsection{The Prolate-Wave Educated Guess $k_\lambda$}

The prolate spheroidal wave operator on $[-\lambda, \lambda]$ is
$\text{PW}_\lambda = -\partial_x((\lambda^2 - x^2)\partial_x)
+ (2\pi\lambda x)^2$.
Let $h_{0,\lambda}$ and $h_{4,\lambda}$ be the two even
eigenfunctions with $0$ and $4$ nodes respectively. Define
$h_\lambda = h_{4,\lambda} - r\,h_{0,\lambda}$ with
$r = \int h_{4,\lambda}\,dx \,/\, \int h_{0,\lambda}\,dx$
so that $\int h_\lambda\,dx = 0$. The educated guess is then
\[
k_\lambda(u) := \sqrt{u}\sum_{n=1}^{\lfloor\lambda/u\rfloor}
h_\lambda(n\,u),
\quad u \in [\lambda^{-1}, \lambda],
\]
i.e.\ the Eisenstein-like sum map applied to $h_\lambda$.

Conjecture~\ref{conj:lemma72} predicts that $k_\lambda$
approximates $\xi_\lambda$ at relative rate $O(\lambda^{-2})$.

\subsection{The Truncated Completed Eta Function $\Lambda_\lambda(s)$}

The truncated completed eta function is
\[
\Lambda_\lambda(s) := \int_{\lambda^{-1}}^{\lambda}
t^{s-1}\,\omega(t)\,dt,
\quad \omega(t) := \frac{t\,e^t}{(1+e^t)^2},
\]
and its $\xi$-weighted variant is
\[
G(s) := \int_{\lambda^{-1}}^{\lambda}
f_\lambda(u)\,\omega(u)\,u^{s-1}\,du,
\]
where $f_\lambda(u) = (1/\sqrt{L})\sum_n \xi_n e^{2\pi i n
\ln(\lambda u)/L}$ is the function reconstructed from
$\xi_\lambda$. As $\lambda \to \infty$, $\Lambda_\lambda(s)
\to \Lambda(s) = \Gamma(s+1)\eta(s)$, whose zeros include the
Riemann zeros~\cite{Yakaboylu2024}.

\section{Implementation}

Our implementation uses the Xcelerator
Toolkit,\footnote{\url{https://github.com/TeamXcelerator/xcelerator-toolkit}}
a publicly available library of high-precision spectral methods for
analytic number theory, pinned to a tagged release for reproducibility.
The toolkit provides the core numerical
primitives (Gauss--Legendre quadrature with disk-cached HP nodes,
LU-based inverse iteration, Newton refinement) and the CCM-specific
construction, prolate-wave operators, and Mellin transforms. The
paper-specific repository contains only the CLI harness and
reproduction scripts.

The toolkit is built on Rust (edition 2021), \texttt{rug} (v1.30,
wrapping MPFR/GMP) for arbitrary-precision arithmetic,
\texttt{nalgebra} (v0.33) for f64 dense linear algebra, and
\texttt{rayon} (v1.10) for parallelism. Three subcommands
correspond to the three tests:
\begin{itemize}
\item \texttt{prolate-compare} (Conjecture~\ref{conj:lemma72}):
loads a pre-computed $\xi_\lambda$ JSON, builds $\text{PW}_\lambda$
via 3-point finite differences on a uniform grid, identifies
$h_{0,\lambda}$ and $h_{4,\lambda}$ by parity and node-counting,
forms $h_\lambda$ and applies the $\mathcal{E}$ map, then computes
$\|\xi_\lambda - c\,k_\lambda\|_\infty / \|\xi_\lambda\|_\infty$
on a logarithmically-spaced sample grid in $[\lambda^{-1}, \lambda]$.
\item \texttt{sliwinski-check} (Conjecture~\ref{conj:sliwinski}):
runs the high-precision CCM construction at given $(\lambda, N)$
configurations, computes the full positive eigenvalue spectrum,
and tabulates $\varepsilon(\lambda, N)$ against
$1/(4\ln\lambda)$ and $\varepsilon\cdot\ln\lambda$.
\item \texttt{mellin-compare} (CCM vs Mellin truncation): scans
the critical line $\Re(s) = 1/2$ for zeros of both $\Lambda_\lambda(s)$
and $G(s)$ via Gauss--Legendre quadrature, then bisects in
high precision to locate each zero precisely.
\end{itemize}

The high-precision $\xi_\lambda$ values are reused from the same
runs that produced the eigenvalue verification in our companion
paper~\cite{Andrews2026A}, ensuring consistency between the two
papers' empirical foundations.

\section{Result 1: CCM Lemma~7.2 is Empirically Falsified}
\label{sec:lemma72}

Conjecture~\ref{conj:lemma72} asserts that
$\text{rel}\times\lambda^2$ should be approximately constant
(equal to $C$) across $\lambda$. We measure this quantity at
three configurations:

\begin{table}[h]
\centering
\caption{Prolate-wave approximation of $\xi_\lambda$ at HP-1000
working precision ($3338$-bit, MPFR/GMP). All three $\xi_\lambda$
values pre-computed at HP-1000 and committed to the data
directory; banded tridiagonal LU recovers each prolate
eigenvector at HP-1000 in milliseconds.}
\label{tab:lemma72}
\begin{tabular}{@{}llrrr@{}}
\toprule
Configuration & $N_{\text{grid}}$ & rel $L^\infty$ error & $\text{rel}\times\lambda^2$ & $\lambda^{-2}$ bound \\
\midrule
$\lambda^2{=}13$, $N{=}120$   & $4001$ & $6.887\times 10^{-5}$ & $8.953\times 10^{-4}$ & $7.69\times 10^{-2}$ \\
$\lambda^2{=}100$, $N{=}500$  & $4001$ & $6.697\times 10^{-4}$ & $6.697\times 10^{-2}$ & $1.00\times 10^{-2}$ \\
$\lambda^2{=}100$, $N{=}500$  & $8001$ & $6.470\times 10^{-4}$ & $6.470\times 10^{-2}$ & $1.00\times 10^{-2}$ \\
$\lambda^2{=}1000$, $N{=}800$ & $8001$ & $\mathbf{9.988\times 10^{-1}}$ & $\mathbf{9.988\times 10^{2}}$ & $1.00\times 10^{-3}$ \\
\bottomrule
\end{tabular}
\end{table}

The growth of $\text{rel}\times\lambda^2$ is dramatic: $8.953\times
10^{-4} \to 6.470\times 10^{-2} \to 9.988\times 10^{2}$ across the
three configurations (using the $N_{\text{grid}}=8001$ row at
$\lambda^2{=}100$ for the apples-to-apples comparison), a factor
of approximately $1.116\times 10^6$. If Conjecture~\ref{conj:lemma72}
held, this column would be approximately constant.

At $\lambda^2 = 1000$, the residual is $99.9\%$ of $\xi_\lambda$
itself, with the optimal scalar~$c$ at exponential magnitude
($\sim 4.567\times 10^{501}$). The educated guess has essentially
no correlation with the actual eigenvector at large~$\lambda$.

\subsection{Grid Convergence Test}

To rule out finite-difference truncation as the cause of the
residual, we doubled the grid resolution at $\lambda^2 = 100$:

\begin{table}[h]
\centering
\caption{Grid convergence of the prolate-wave residual at
$\lambda^2{=}100$, HP-1000. Halving the FD grid spacing changes
the residual by $3.6\%$, far short of the $4\times$ reduction
an $O(h^2)$ method would produce if FD truncation dominated.}
\label{tab:grid_conv}
\begin{tabular}{@{}lrl@{}}
\toprule
$N_{\text{grid}}$ & rel $L^\infty$ & change from $4001$ \\
\midrule
$4001$ & $6.71253\times 10^{-4}$ & --- \\
$8001$ & $6.46957\times 10^{-4}$ & $-3.6\%$ \\
\bottomrule
\end{tabular}
\end{table}

If the residual were dominated by FD truncation, doubling the
grid (halving $h$) should cut the error by $4\times$ (an
$O(h^2)$ method). Instead it changed by $3.6\%$, ruling out
truncation as the driver.

\subsection{Interpretation}

The prolate operator itself is well-resolved: the smallest
prolate eigenvalue at $\lambda^2 = 1000$ is $6282.28$, against the
asymptotic prediction $2\pi\lambda^2 = 6283.19$ ($0.014\%$ relative
error). The failure is not in the prolate eigenfunctions, but in
the chain $h_{0,\lambda}, h_{4,\lambda} \to h_\lambda
\to k_\lambda \to$ comparison with $\xi_\lambda$. The educated
guess works well at small $\lambda$ (where $\text{PW}_\lambda$
approaches the harmonic oscillator limit) and breaks down at
large~$\lambda$.

\begin{remark}
This finding is consistent with the post-2022 CCM literature
trajectory: the authors moved away from the positive prolate
spectrum in their 2022 PNAS work (toward the Sonin space), and
\'Sliwi\'nski's January 2026 paper proves an inverse-logarithmic
lower bound via Heisenberg uncertainty that does not rely on
prolate waves at all.
\end{remark}

\section{Result 2: \'Sliwi\'nski's Conjecture~4.1 is Empirically Unsupported at HP-1000}
\label{sec:sliwinski}

\'Sliwi\'nski~\cite{Sliwinski2026} proves Theorem~3.1, the lower
bound $\varepsilon(\lambda, N) \geq 1/(4\ln\lambda)$, via a
Heisenberg uncertainty argument: position operator on
$L^2([\lambda^{-1}, \lambda], du)$ has support size $2\ln\lambda$,
hence $\sigma_u \leq 2\ln\lambda$, and $\sigma_u \cdot \sigma_D
\geq 1/2$ gives $\sigma_D \geq 1/(4\ln\lambda)$.
Conjecture~\ref{conj:sliwinski} further claims this bound is
asymptotically tight: $\varepsilon(\kappa)\ln(\kappa)$ approaches
a nonzero constant.

\subsection{Standard CCM Regime ($N \gg \lambda$)}

\begin{table}[h]
\centering
\caption{\'Sliwi\'nski's lower bound applied to our standard CCM
configurations, all at HP-1000 ($3338$-bit) working precision.
$\varepsilon$ is the mean absolute spectral error against the
first $N$ Riemann zeros.}
\label{tab:sliw_standard}
\begin{tabular}{@{}rrrrrrrr@{}}
\toprule
$\lambda$ & $\lambda^2$ & $N$ & digits & mean abs $\varepsilon$ & $1/(4\ln\lambda)$ & ratio & $\varepsilon\cdot\ln\lambda$ \\
\midrule
$\sqrt{13}$       & $13$    & $120$ & $1000$ & $1.063$  & $0.195$  & $5.454\times$  & $1.363$ \\
$10$              & $100$   & $500$ & $1000$ & $1.595$  & $0.109$  & $14.69\times$  & $3.673$ \\
$\sqrt{1000}$     & $1000$  & $800$ & $1000$ & $3.608$\textsuperscript{\dag} & $0.072$  & $49.85\times$ & $12.46$ \\
\bottomrule
\multicolumn{8}{l}{\footnotesize\textsuperscript{\dag}At HP-500 this
row previously read $\varepsilon = 5.638$ ($\varepsilon\cdot\ln\lambda
= 19.5$). That value was inflated by} \\
\multicolumn{8}{l}{\footnotesize working-precision saturation of the
low-index eigenvalues: at HP-500 the first eigenvalues saturate near
$500$ matching} \\
\multicolumn{8}{l}{\footnotesize digits, whereas at HP-1000 they
resolve to ${\sim}999$ (Table~\ref{tab:per_eigenvalue}), lowering the
mean. The HP-1000 value is the accurate one.} \\
\end{tabular}
\end{table}

Theorem~3.1 holds with substantial margin ($5\times$ to $50\times$).
The conjecture, however, would require $\varepsilon\cdot\ln\lambda$
to be roughly constant; we observe it growing from $1.36$ to
$12.46$ across the three configurations (a factor of $9.14$).

\subsection{\'Sliwi\'nski's Own Regime ($\kappa = N = \lambda$) at HP-1000}
\label{subsec:sliw_hp1000}

This is the regime in which Conjecture~\ref{conj:sliwinski} was
formulated. We test it at $1000$-digit working precision
($3338$~bits, MPFR/GMP) across $\kappa \in \{50, 100, 150, 200,
300, 400, 500\}$. The full-spectrum metric $\bar\varepsilon = (1/N)
\sum_k |\nu_k - \zeta_k|$ uses all $N$ eigenvalues; the
trimmed-interior metric drops the top $20\%$ of edge
eigenvalues, which lie at the matrix's truncation boundary
(\emph{cf.}\ Section~\ref{subsec:per_eigvalue}).

\begin{table}[h]
\centering
\caption{\'Sliwi\'nski's regime $\kappa = N = \lambda$ at HP-1000
($3338$-bit precision). Full-spectrum metric uses all $N$
eigenvalues; trim-$20\%$ drops the $\lfloor 0.2 N\rfloor$
highest-index eigenvalues. ``ratio'' is the Theorem~3.1 ratio
$\varepsilon / [1/(4\ln\lambda)]$; values $> 1$ mean the bound
is satisfied.}
\label{tab:sliw_hp1000}
\begin{tabular}{@{}rrrrrrr@{}}
\toprule
$\kappa$ & full $\bar\varepsilon$ & full $\varepsilon{\cdot}\ln\lambda$ & full ratio & trim $\bar\varepsilon$ & trim $\varepsilon{\cdot}\ln\lambda$ & trim ratio \\
\midrule
$50$  & $2.938$ & $11.49$  & $45.97\times$  & $0.3196$ & $1.250$  & $5.001\times$  \\
$100$ & $1.692$ & $7.792$  & $31.17\times$  & $0.9644$ & $4.441$  & $17.77\times$ \\
$150$ & $2.984$ & $14.95$  & $59.82\times$  & $0.5834$ & $2.923$  & $11.69\times$ \\
$200$ & $1.579$ & $8.367$  & $33.47\times$  & $0.9123$ & $4.834$  & $19.34\times$ \\
$300$ & $2.156$ & $12.30$  & $49.20\times$  & $0.8236$ & $4.698$  & $18.79\times$ \\
$400$ & $1.310$ & $7.846$  & $31.38\times$  & $0.7553$ & $4.525$  & $18.10\times$ \\
$500$ & $5.163$ & $\mathbf{32.09}$ & $128.4\times$ & $1.699$ & $\mathbf{10.56}$ & $42.23\times$ \\
\bottomrule
\end{tabular}
\end{table}

Two findings:

\begin{enumerate}
\item \textbf{Theorem~3.1 holds at every configuration with
substantial margin}, in both metrics: ratios from $5.0\times$
($\kappa{=}50$ trim) to $128.4\times$ ($\kappa{=}500$ full).
The bound is rigorously satisfied across the entire HP-1000 sweep.

\item \textbf{Conjecture~4.1 is empirically unsupported at HP-1000.}
The conjecture predicts $\varepsilon(\kappa)\cdot\ln\kappa$
approaches a single nonzero constant. The data does not approach
a constant in the tested range. Across $\kappa \in [50, 500]$:

\begin{itemize}
\item Full-spectrum: $\varepsilon\cdot\ln\lambda$ ranges over
$[7.792, 32.09]$, a $4.118\times$ spread, with first/last
($\kappa{=}50 \to 500$) ratio $2.792\times$.
\item Trim-$20\%$: ranges over $[1.250, 10.56]$, a $8.445\times$
first/last ratio (both extrema are the endpoints).
\item Both metrics show a sharp upturn at $\kappa = 500$: $4.090\times$
jump in the full metric over $\kappa{=}400$, $2.333\times$ in the
trimmed metric.
\end{itemize}
\end{enumerate}

The $\kappa = 500$ upturn may reflect a structural matrix
resonance at this specific $(\kappa, N)$ pair, or numerical
conditioning in the inverse-iteration and Newton refinement
stages near the matrix's truncation boundary; we report both
possibilities. Either way, $\varepsilon(\kappa)\cdot\ln\kappa$
does not approach a constant in the tested range, and the
conjectural limit form is not supported by this data.

\paragraph{Floor independence.}
At $\kappa = 400$ and $500$ the first eigenvalues saturate at the
HP-1000 precision ceiling ($\sim 999$ matching digits; see
Table~\ref{tab:per_eigenvalue}), meaning $\varepsilon_N$ at these
configs is at or below the working-precision floor. This does
\emph{not} invalidate the aggregate metric $\varepsilon(\kappa)$:
the mean absolute error is dominated by the high-index eigenvalues
at the matrix's truncation boundary, which contribute $O(1)$ error
independently of the precision floor. Raising the working precision
to HP-2000 would sharpen the first few eigenvalues (from $999$ to
${\sim}2000$ digits) but would not materially change
$\varepsilon(\kappa)$, because the $O(1)$ truncation-boundary tail
still dominates the sum. The upturn at $\kappa = 500$ is therefore
structural (more eigenvalues at the boundary contributing larger
errors), not a precision-floor artifact.

The companion repository contains the trimmed-metric flag
(\texttt{--trim-edge-fraction}) used to produce the
right-hand columns above, so reviewers can re-run the analysis
at any precision and any trim fraction they choose.

\subsection{Mean vs.\ Uniform Error}

\'Sliwi\'nski's Conjecture~4.1 targets the uniform (pointwise
maximum) error $\max_k |\nu_k - \zeta_k|$ rather than the mean
$(1/N)\sum_k |\nu_k - \zeta_k|$ used as the primary metric in
Section~\ref{sec:sliwinski}; we thank \'Sliwi\'nski (private
communication, 2026) for clarifying this distinction. Our data
shows the uniform error also fails to stabilize when multiplied
by $\ln^\alpha(\kappa)$ for any $\alpha$:

\begin{table}[h]
\centering
\caption{Uniform error $\times\,\ln^\alpha(\kappa)$ at HP-1000
($3338$-bit precision). If the uniform error decayed as
$1/\ln^\alpha$, one column would be approximately constant. None
do --- the values oscillate over orders of magnitude across
$\kappa$, with a sharp upturn at $\kappa = 500$.}
\label{tab:uniform_alpha}
\begin{tabular}{@{}rrrrr@{}}
\toprule
$\kappa$ & uniform err & $\text{unif}\cdot\ln^1$ & $\text{unif}\cdot\ln^2$ & $\text{unif}\cdot\ln^3$ \\
\midrule
$50$  & $66.05$ & $258.4$ & $1011$   & $3955$   \\
$100$ & $38.87$ & $179.0$ & $824.3$  & $3796$   \\
$150$ & $154.5$ & $774.3$ & $3880$   & $19440$  \\
$200$ & $60.59$ & $321.0$ & $1701$   & $9013$   \\
$300$ & $139.5$ & $795.8$ & $4539$   & $25890$  \\
$400$ & $52.70$ & $315.7$ & $1892$   & $11340$  \\
$500$ & $629.1$ & $3910$  & $24300$  & $151000$ \\
\bottomrule
\end{tabular}
\end{table}

No column stabilizes. The uniform error is dominated by the
\emph{last} eigenvalue ($k = N$), which stays at $\lesssim 3$
matching digits regardless of $\kappa$ (\emph{cf.}\
Table~\ref{tab:per_eigenvalue}, last column).

\'Sliwi\'nski (private communication) further suggested a
softer extended-form conjecture --- $1/\ln^\alpha(\kappa)$ for
some $\alpha > 1$, rather than the original $\alpha = 1$.
Table~\ref{tab:uniform_alpha} tests this directly: if the
uniform error decayed as $1/\ln^\alpha(\kappa)$ for any fixed
$\alpha \in \{1, 2, 3\}$, the corresponding column would be
approximately constant across $\kappa$. None of the columns
$\alpha \in \{1, 2, 3\}$ stabilize at HP-1000 in the tested
range $\kappa \in [50, 500]$; the values vary by factors of
$\sim 16\times$ ($\alpha = 0$, raw uniform) to $\sim 40\times$
($\alpha = 3$) within each column, and
the sharp upturn at $\kappa = 500$ ($\sim 12\times$ over
$\kappa = 400$) is preserved in every column independent of
$\alpha$. The extended-form conjecture is therefore also
unsupported at HP-1000 in this regime, with the same caveat
as before: the conjecture is asymptotic and our test is finite.

\subsection{Per-Eigenvalue Convergence Profile at HP-1000}
\label{subsec:per_eigvalue}

The HP-1000 spectrum exhibits a sharp split between converged
and unconverged eigenvalues:

\begin{table}[h]
\centering
\caption{Matching digits per eigenvalue index at HP-1000. Each
entry is $-\log_{10}|\nu_k - \zeta_k|/|\zeta_k|$ computed in HP
arithmetic. The first eigenvalues converge super-exponentially;
the last eigenvalue ($k = N$) remains at $\lesssim 3$ matching
digits regardless of $\kappa$.}
\label{tab:per_eigenvalue}
\begin{tabular}{@{}rrrrrr@{}}
\toprule
$\kappa$ & $k{=}1$ & $k{=}5$ & $k{=}10$ & $k{=}20$ & $k{=}N$ \\
\midrule
$50$  & $160.8$ & $126.5$ & $61.2$  & $15.5$  & $2.2$ \\
$100$ & $315.1$ & $290.3$ & $254.0$ & $135.7$ & $2.8$ \\
$150$ & $463.3$ & $441.8$ & $412.8$ & $340.7$ & $2.6$ \\
$200$ & $608.4$ & $588.9$ & $563.3$ & $504.9$ & $0.8$ \\
$300$ & $893.0$ & $875.7$ & $853.9$ & $806.9$ & $1.3$ \\
$400$ & $999.4$ & $1001$  & $1000$  & $988.8$ & $2.4$ \\
$500$ & $999.4$ & $1001$  & $1000$  & $1000$  & $1.9$ \\
\bottomrule
\end{tabular}
\end{table}

The first eigenvalue's matching-digit count grows roughly
linearly with $\kappa$ through $\kappa = 300$ (160 → 315 → 463
→ 608 → 893, an additive trend of $\sim$140--150 digits per
$\Delta\kappa = 50$), then saturates against the $1000$-digit
working-precision ceiling at $\kappa = 400$ and beyond
($999 \to 999$). At $\kappa = 200$ the first $20$ eigenvalues
all match to $> 500$ digits; at $\kappa = 300$, $> 800$; at
$\kappa = 400$, $> 988$; at $\kappa = 500$, all of the first
$20$ are at the precision ceiling. The last eigenvalue
($k = N$) remains at $\lesssim 3$ matching digits regardless of
$\kappa$, confirming that the aggregate metrics in
Table~\ref{tab:sliw_hp1000} are dominated by the unconverged
tail at the truncation boundary, not by any intrinsic
convergence rate of the construction.

\begin{remark}
The $\sqrt{N}$ growth of mean error visible in
Table~\ref{tab:sliw_hp1000} (full-spectrum mean rising from
$1.69$ at $\kappa{=}100$ to $5.16$ at $\kappa{=}500$, roughly
$\sqrt{5}\times$ for a $5\times$ increase in $N$) is consistent
with a central limit theorem effect: if the first $O(1)$
eigenvalues converge to full precision and the remaining $N - O(1)$
contribute $O(1)$ error each, the mean over $N$ terms grows
as $\sqrt{N}$ by the CLT applied to the GUE-distributed zero
spacings. This connection was noted by
\'Sliwi\'nski~(private communication, 2026).
\end{remark}

\section{Result 3: CCM is $10^{55}$ More Accurate than Naive Mellin Truncation}
\label{sec:mellin}

If a substantial fraction of CCM's accuracy could be captured by
an integral transform of $\xi_\lambda$, that would constitute a
candidate bridge to RH. We test this directly.

\subsection{Naive Mellin Truncation}

We computed zeros of $\Lambda_\lambda(s)$ on the critical line
$\Re(s) = 1/2$ via Gauss--Legendre quadrature with $500$~points
at HP-1000 precision ($3338$~bits), followed by HP bisection.

\begin{table}[h]
\centering
\caption{Zeros of the truncated completed eta function $\Lambda_\lambda(s)$
on the critical line, $\lambda = \sqrt{13}$, HP-1000.}
\label{tab:lambda_sqrt13}
\begin{tabular}{@{}rrrr@{}}
\toprule
$k$ & $\Lambda_\lambda$ zero & Riemann zero $\zeta_k$ & difference \\
\midrule
$1$ & $14.797$ & $14.135$ & $0.6622$ \\
$2$ & $22.112$ & $21.022$ & $1.090$ \\
$3$ & $24.555$ & $25.011$ & $0.4557$ \\
$4$ & $29.445$ & $30.425$ & $0.9803$ \\
$5$ & $31.890$ & $32.935$ & $1.045$ \\
\bottomrule
\end{tabular}
\end{table}

\begin{table}[h]
\centering
\caption{Same comparison at $\lambda = 10$ ($\lambda^2 = 100$).}
\label{tab:lambda_10}
\begin{tabular}{@{}rrrr@{}}
\toprule
$k$ & $\Lambda_\lambda$ zero & Riemann zero $\zeta_k$ & difference \\
\midrule
$1$ & $13.719$ & $14.135$ & $0.416$ \\
$2$ & $20.518$ & $21.022$ & $0.504$ \\
$3$ & $24.603$ & $25.011$ & $0.408$ \\
$4$ & $30.053$ & $30.425$ & $0.372$ \\
$5$ & $32.779$ & $32.935$ & $0.156$ \\
\bottomrule
\end{tabular}
\end{table}

In both cases the absolute error is $O(0.1\text{--}1.0)$, i.e.
\emph{less than one matching digit}. The CCM construction at the
same $(\lambda, N)$ matches Riemann zeros to $55$~digits at
$\lambda^2 = 13$, $460$~digits at $\lambda^2 = 100$, and $999$~digits
at $\lambda^2 = 1000$ \cite{Andrews2026A}. The improvement factor
is therefore $\sim 10^{55}$, $\sim 10^{460}$, and $\sim 10^{999}$
across the three configurations.

\subsection{$\xi$-Weighted Mellin Transform}

If $\xi_\lambda$ encodes essential information about Riemann
zeros, weighting the Mellin integrand by $f_\lambda(u)$ might
recover some of the CCM accuracy. We test this with the
$\xi$-weighted variant $G(s)$ at $\lambda^2 = 13$ at HP-1000
($3338$-bit) precision over the scan range $\Im(s)\in[5,55]$:

\begin{table}[h]
\centering
\caption{$\xi$-weighted Mellin $G(s)$ vs naive $\Lambda_\lambda$ at
$\lambda = \sqrt{13}$, $N = 120$, full HP-1000 arithmetic
($3338$-bit precision, $500$-point Gauss--Legendre quadrature,
$5000$ scan points). Eleven Riemann zeros lie in the scan range.}
\label{tab:xi_weighted}
\setlength{\tabcolsep}{4pt}
\small
\begin{tabular}{@{}rrrrrrr@{}}
\toprule
$k$ & $\Lambda_\lambda$ zero & $G(s)$ zero & Riemann zero & $\Lambda_\lambda$ error & $G$ error & improvement \\
\midrule
$1$  & $14.797$ & $13.930$ & $14.135$ & $0.6622$ & $0.2045$ & $3.238\times$  \\
$2$  & $22.112$ & $20.815$ & $21.022$ & $1.090$  & $0.2067$ & $5.275\times$  \\
$3$  & $24.555$ & $25.152$ & $25.011$ & $0.4557$ & $0.1411$ & $3.230\times$  \\
$4$  & $29.445$ & $30.336$ & $30.425$ & $0.9803$ & $0.08852$ & $11.07\times$  \\
$5$  & $31.890$ & $33.399$ & $32.935$ & $1.045$  & $0.4636$ & $2.253\times$  \\
$6$  & $36.784$ & $37.725$ & $37.586$ & $0.8026$ & $0.1389$ & $5.779\times$  \\
$7$  & $41.678$ & $41.322$ & $40.919$ & $0.7595$ & $0.4034$ & $1.883\times$  \\
$8$  & $44.126$ & $43.829$ & $43.327$ & $0.7989$ & $0.5018$ & $1.592\times$  \\
$9$  & $49.022$ & $48.162$ & $48.005$ & $1.017$  & $0.1570$ & $6.479\times$  \\
$10$ & $49.022$ & $50.585$ & $49.774$ & $0.7519$ & $0.8107$ & $0.9274\times$ \\
$11$ & $53.919$ & $53.480$ & $52.970$ & $0.9482$ & $0.4384$ & $2.163\times$  \\
\bottomrule
\end{tabular}
\end{table}

Mean $G$ error: $0.3231$. Mean $\Lambda_\lambda$ error: $0.8464$.
Mean improvement factor: $\mathbf{2.619\times}$.

\begin{remark}[Zero counts on the critical line]
\label{rem:zero_counts}
The $\Lambda_\lambda$ scan over $\Im(s)\in[5,55]$ at $\lambda^2=13$
finds $21$ zeros, while the $\xi$-weighted $G(s)$ scan over the
same range finds $11$ zeros --- exactly the count of Riemann zeros
in the range. The unweighted truncation
introduces approximately ten spurious zeros that are not present
in the weighted variant. This is consistent with $\xi$-weighting
suppressing oscillatory artefacts of the bare integrand $\omega(t)$
near the truncation endpoints, while the $11$ zeros that survive
align (modulo the residual $O(0.1\text{--}0.5)$ errors above) with
the actual Riemann zeros.
\end{remark}

The $\xi$-weighting helps modestly but does not approach the
$10^{55}$ accuracy of the full CCM construction. The
HP-1000 measurement reproduces the previously reported
$2.619\times$ improvement to four significant figures, confirming
that the residual $O(0.1\text{--}0.5)$ errors are mathematical in
nature, not numerical artefacts: at higher precision the same
answer emerges.

\subsection{Interpretation}

The $10^{55}$ improvement factor of CCM over Mellin truncation
quantifies what the Weil quadratic form's algebraic structure
contributes \emph{beyond} integration:
\begin{itemize}
\item The eigenvalue problem $\tau\xi = \varepsilon_N\xi$ extracts
the smallest eigenpair with full precision available.
\item The rational function $R(z) = \xi_0 + 2z\sum_j \xi_j/(z - j^2)$
extracts zeros via an algebraic mechanism whose accuracy is bounded
only by $\varepsilon_N$.
\item Neither operation has an integral-transform analogue that
captures comparable information from $\xi_\lambda$ alone.
\end{itemize}

The implication for the bridge problem is sharp: no integral
transform of $\xi_\lambda$ tested here approaches the accuracy
CCM exhibits, and the structural argument above (eigenvalues
together with rational-function zeros lose information under
integration) suggests this generalizes. The bridge must engage
with the algebraic mechanism (eigenvalue problem followed by
rational-function zeros).

\section{Discussion}

\subsection{Status of the Two Quantitative Bridges}

Conjecture~\ref{conj:lemma72} proposed an analytic bridge from
$\xi_\lambda$ to the Riemann zeros via prolate-wave approximation
at rate $O(\lambda^{-2})$. Our high-precision data falsifies the
predicted rate: $\text{rel}\times\lambda^2$ grows by $\sim
1.116 \times 10^6$ across $\lambda^2 \in [13, 1000]$. At $\lambda^2 =
1000$ the residual is $99.9\%$ of $\xi_\lambda$ itself, with the
optimal scalar $c$ at exponential magnitude ($\sim 4.567\times
10^{501}$); the educated guess has essentially no correlation
with the actual eigenvector at large $\lambda$.

Conjecture~\ref{conj:sliwinski} proposed an asymptotic rate
$\varepsilon \sim 1/\ln\lambda$ in the regime $\kappa = N =
\lambda$. At HP-1000 ($3338$-bit) across $\kappa \in [50, 500]$
(Section~\ref{subsec:sliw_hp1000}), $\varepsilon\cdot\ln\lambda$
ranges over $[7.792, 32.09]$ (full-spectrum, $4.118\times$ spread)
and $[1.250, 10.56]$ (trimmed-interior, $8.445\times$ first/last),
with a sharp upturn at $\kappa = 500$ in both metrics. Neither
trajectory is consistent with the conjectural limit
$\varepsilon(\kappa)\cdot\ln(\kappa) \to C$. Theorem~3.1
(\'Sliwi\'nski's lower bound) holds with substantial margin at
all configurations ($5.001\times$--$128.4\times$ ratios). Whether the
$\kappa = 500$ upturn reflects a structural matrix resonance or
numerical conditioning at the matrix's truncation boundary is
under investigation; either way, the data does not support the
predicted asymptotic constancy in the tested range.

\subsection{The Algebraic Picture}

The per-eigenvalue analysis (Table~\ref{tab:per_eigenvalue})
reveals that the CCM construction does not have a convergence
rate problem for the first eigenvalues. In the \'Sliwi\'nski
regime ($\kappa = N = \lambda$), the first eigenvalue matches
the first Riemann zero to \emph{$160.8$~digits at $\kappa = 50$
and $608.4$~digits at $\kappa = 200$}, with the matching-digit
count growing roughly linearly in $\kappa$ through $\kappa =
300$ (additive trend $\sim$140--150 digits per $\Delta\kappa =
50$) before saturating against the $1000$-digit working-precision
ceiling at $\kappa = 400$ and beyond. At $\kappa = 200$, the
first $20$ eigenvalues all match Riemann zeros to ${>}500$
digits; at $\kappa = 300$, ${>}800$; at $\kappa = 400$ and
beyond, the first $20$ are at the precision ceiling.

The construction's accuracy on these eigenvalues is limited
only by arithmetic precision, not by any intrinsic convergence
rate. No rate of the form $1/\ln^\alpha(\kappa)$ describes this
behavior --- it is super-exponential convergence that immediately
saturates at whatever precision is available.

The ``slow convergence'' measured by the mean and uniform errors
comes entirely from the \emph{last} eigenvalue ($k = N$), which
remains at ${\sim}1$--$3$ matching digits regardless of $\kappa$
even at HP-1000 working precision. This is a structural property
of the matrix's truncation boundary, not of the working precision.
The correct question is not ``what is the convergence rate?'' but
rather ``how many eigenvalues converge to full precision as
$\kappa$ grows?'' The HP-1000 per-eigenvalue table shows this
count growing rapidly with $\kappa$; characterizing this growth
precisely --- and proving it unbounded --- remains open.

Combined with the $10^{55}$ improvement of CCM over naive Mellin
truncation, the picture that emerges is:

\begin{itemize}
\item The CCM construction's accuracy is \emph{algebraic}, not
analytic. It comes from solving an eigenvalue problem on a
specific symmetric matrix, then finding zeros of a rational
function determined by the eigenvector.
\item Neither step has an integral-transform analogue that
preserves the full information. Hence no Mellin-style bridge can
recover comparable accuracy.
\item The candidates studied in the literature (prolate
approximation, inverse-logarithmic rate, naive or $\xi$-weighted
Mellin) all fall short by orders of magnitude when tested at
high precision.
\end{itemize}

The construction works (Step~1 of the CCM proof strategy is
verified to $999$~digits at $\lambda^2{=}1000$, $N{=}800$, HP-1000
in~\cite{Andrews2026A}) but the
quantitative bridges proposed so far do not. A successful
proof of RH via this route requires a new framework that engages
with the algebraic structure directly and predicts a rate
consistent with high-precision data.

\subsection{What Remains}

Two directions remain open:
\begin{enumerate}
\item A different rate prediction in the $\kappa = N = \lambda$
regime, consistent with our $\varepsilon \sim \sqrt{N}$
observation, would refine \'Sliwi\'nski's framework.
\item A new bridge construction that does \emph{not} go through
prolate waves or Mellin transforms could complete the CCM proof
strategy. Candidate frameworks in the literature (CCM 2022 PNAS
Sonin space, Yakaboylu's Hilbert--P\'olya Hamiltonian) share
operators with CCM 2025 but do not yet provide explicit bridges
from~$\xi_\lambda$.
\end{enumerate}

\section{Conclusion}

We have empirically tested two published quantitative predictions
about the CCM zeta spectral triple's convergence rate at high
precision. CCM Lemma~7.2 fails by a factor of approximately
$1.116\times 10^6$ in $\text{rel}\times\lambda^2$ across $\lambda^2
\in [13, 1000]$. \'Sliwi\'nski's Conjecture~4.1 is empirically
unsupported at HP-1000: across $\kappa = N = \lambda \in [50, 500]$,
$\varepsilon\cdot\ln\lambda$ ranges over $[7.792, 32.09]$
(full-spectrum) and $[1.250, 10.56]$ (trimmed-interior, an
$8.445\times$ first/last change), with a sharp upturn at
$\kappa = 500$, neither metric approaching a single nonzero
constant. Theorem~3.1 (\'Sliwi\'nski's lower bound) holds with
substantial margin throughout. The CCM construction is
approximately $10^{55}$ more accurate than naive Mellin
truncation, and $\xi$-weighting the Mellin integrand recovers
only a $2.6\times$ improvement (verified at HP-1000).

The per-eigenvalue analysis at HP-1000 reveals a more nuanced
picture than simple falsification. The first eigenvalues converge
\emph{super-exponentially} --- the first Riemann zero is matched
to $160.8$~digits at $\kappa = 50$ and to $608.4$~digits at
$\kappa = 200$, with linear growth in $\kappa$ through $\kappa
= 300$ before saturating against the $1000$-digit
working-precision ceiling at $\kappa = 400$. The aggregate error metrics (mean,
uniform) are dominated by high-index eigenvalues lying at the
matrix's truncation boundary. Even at HP-1000 the last
eigenvalue ($k = N$) remains at $\sim 1$--$3$ matching digits
regardless of $\kappa$, indicating a structural limit at the
boundary rather than a numerical one. The $\sqrt{N}$ growth of
the mean error visible in our data is consistent with a CLT
effect on GUE-distributed zero spacings.

The construction itself is not in question --- it works
extraordinarily well on the first eigenvalues. What is in
question is the route from $\xi_\lambda$ to the Riemann
Hypothesis. The correct convergence question is: how does the
number of fully-converged eigenvalues grow with $\kappa$? Our
HP-1000 data shows the first ${\sim}20$ eigenvalues all matching
Riemann zeros to ${>}500$ digits at $\kappa = 200$, with the
count of fully-converged eigenvalues growing rapidly with
$\kappa$; characterizing this growth precisely --- and proving
it unbounded --- remains open.

\section*{Acknowledgments}

All computations were performed on Vast.ai cloud infrastructure
(AMD EPYC class servers, varying core counts across runs); the
author personally funded all compute resources. The author identified the research questions, set the
research direction, and made all decisions about what experiments
to run and how to interpret the results. Specific contributions
include: the convergence question that motivated this empirical
testing of CCM and \'Sliwi\'nski's predictions; the strategy of
parallelizing high-precision computations to make 200--500 digit
sweeps feasible at scale; the recognition that Gauss--Legendre
nodes and weights are uniquely determined by precision and node
count and thus reusable across runs (yielding a substantial
speedup); and the decision to compute at high precision throughout
to remove numerical noise as a confounding factor in the
falsification claims. AI coding tools (Claude, Anthropic) were
used for Rust implementation under the author's direction; all
research decisions, interpretations, and intellectual
contributions are the author's. The high-precision $\xi_\lambda$
values used in the prolate-wave and Mellin tests come from the
same multi-hour HP runs that produced the eigenvalue verification
in our companion paper~\cite{Andrews2026A}.

\section*{Data and Code Availability}

The complete source code, reference data, and reproduction
scripts are available at:\\
{\small\url{https://github.com/TeamXcelerator/ccm-convergence-rate-falsifications}}

\noindent The shared mathematical library (Xcelerator Toolkit) is at:\\
{\small\url{https://github.com/TeamXcelerator/xcelerator-toolkit}}

\begin{thebibliography}{9}

\bibitem{CCM2025}
A.~Connes, C.~Consani, H.~Moscovici,
\emph{Zeta Spectral Triples},
arXiv:2511.22755 (2025).

\bibitem{Sliwinski2026}
D.~\'Sliwi\'nski,
\emph{Spectral Analysis of the $D_{\log}^{(\lambda, N)}$ Operators},
arXiv:2601.12133 (2026).

\bibitem{Andrews2026A}
R.~Andrews, Jr.,
\emph{Independent Reproduction and Convergence Analysis of the
Connes--Consani--Moscovici Zeta Spectral Triple}.
GitHub: TeamXcelerator/ccm-reproduction-and-convergence (2026).

\bibitem{Yakaboylu2024}
E.~Yakaboylu,
\emph{Hamiltonian for the Hilbert--P\'olya Conjecture},
J.\ Phys.\ A: Math.\ Theor.\ 57, 235204 (2024).
arXiv:2309.00405; see also arXiv:2408.15135.

\bibitem{Odlyzko}
A.~M.~Odlyzko,
\emph{Tables of zeros of the Riemann zeta function},
\url{http://www.dtc.umn.edu/~odlyzko/zeta_tables/}.

\bibitem{PARI}
The PARI Group,
\emph{PARI/GP version 2.15},
Universit\'e de Bordeaux, 2023.
\url{http://pari.math.u-bordeaux.fr/}.

\end{thebibliography}

\end{document}
